% Generated draft: requires mathematical and source review.
\documentclass[11pt]{article}
\usepackage[margin=25mm]{geometry}
\usepackage{fontspec}
\setmainfont{Latin Modern Roman}
\usepackage{amsmath,amssymb,mathtools,unicode-math}
\setmathfont{Latin Modern Math}
\usepackage{enumitem,needspace,xurl,hyperref}
\setlength{\parindent}{0pt}
\setlength{\parskip}{5pt}
\setlength{\emergencystretch}{3em}
\newcommand{\N}{\mathbb{N}}
\newcommand{\Z}{\mathbb{Z}}
\newcommand{\Q}{\mathbb{Q}}
\newcommand{\R}{\mathbb{R}}
\newcommand{\C}{\mathbb{C}}
\newcommand{\F}{\mathbb{F}}
\newcommand{\A}{\mathbb{A}}
\newcommand{\PP}{\mathbb{P}}
\newcommand{\E}{\mathbb{E}}
\newcommand{\eps}{\varepsilon}
\newcommand{\dd}{\,\mathrm{d}}
\begin{document}
\section*{UnsolvedMath ID: 701485}
\textbf{Source problem code:} catalogue entry
\subsection{Definitions and mathematical statement}
A finite \(2\)-group is a group \(G\) with \(|G|=2^n\) for some integer \(n\geq 0\). For a subgroup \(H\leq G\), define its core in \(G\) by
\[
H_G=\bigcap_{g\in G}gHg^{-1}.
\]
This is the largest normal subgroup of \(G\) contained in \(H\), and \(|H/H_G|=[H:H_G]\).

The question proposed by G. Cutolo is whether, for every finite \(2\)-group \(G\) with \(|G|>2\), the condition
\[
\forall H\leq G,\qquad [H:H_G]\leq 2
\]
implies that there exists a subgroup \(A\leq G\) such that
\[
[G:A]=4
\qquad\text{and}\qquad
\forall a,b\in A,\quad ab=ba.
\]
The index required is exactly \(4\); the trivial subgroup is permitted. The supplied registry labels this question open; that label is not an independently verified assessment.
\subsection{Short English statement}
Must every finite \(2\)-group of order greater than \(2\), in which each subgroup has index at most \(2\) over its core, contain an abelian subgroup of index exactly \(4\)?
\subsection{Sources}
\begin{itemize}
\item Local UnsolvedMath registry, supplied complete source record: permanent ID \(701485\), exact source code \texttt{catalogue entry}, title ``Kourovka Notebook Problem 18.56.'' The supplied status is \texttt{open}.
\item \emph{The Kourovka Notebook (Unsolved Problems in Group Theory)}, problem 18.56, attributed to G. Cutolo, proposed in 2014. The supplied background identifies the 18th issue (2014) and the 21st issue edition (2026), with arXiv identifier \texttt{1401.0300v48}. The supplied collection description names E. I. Khukhro and V. D. Mazurov as editors. These bibliographic details are supplied attribution and have not been independently verified.
\end{itemize}
\end{document}
